% Source: Evans, "Partial Differential Equations" (2nd ed., AMS Graduate Studies in Mathematics vol. 19, 2010),
% Chapter 2 "Four Important Linear PDE", PDF pages 27-99 of the cited source.
% Transcribed from the cited source and organized according to leanblueprint conventions.

\section{Transport Equation}
\label{sec:transport-equation}

Fix $b\in\R^n$. The constant-coefficient transport equation for
$u:\R^n\times[0,\infty)\to\R$ is
$$u_t+b\cdot Du=0,$$
where $D=D_x$. Its characteristic curves have direction $(b,1)$.

\subsection{Initial-Value Problem}
\label{sec:transport-equation-initial-value-problem}

Consider
$$(2)\qquad
\begin{cases}
u_t+b\cdot Du=0&\text{in }\R^n\times(0,\infty),\\
u=g&\text{on }\R^n\times\{t=0\}.
\end{cases}$$

\begin{definition}[Explicit solution of the transport equation]
\label{def:transport-explicit-solution}
\lean{EvansLib.transportSolution}
\leanok
\dcref{ch2:2.1.1-def-1}
\group{Transport Equation}


For $b\in\R^n$ and $g:\R^n\to\R$, define
$$u(x,t):=g(x-tb)\qquad(x\in\R^n,\ t\in\R).$$
\end{definition}

\begin{theorem}[Solution of the transport initial-value problem]
\label{thm:transport-ivp-solution}
\lean{EvansLib.transportSolution_isClassicalSolution}
\uses{def:transport-explicit-solution}
\leanok
\dcref{ch2:2.1.1-thm-1}
\group{Transport Equation}



If $g\in C^1(\R^n)$, then $u(x,t)=g(x-tb)$ belongs to
$C^1(\R^n\times\R)$ and satisfies $u_t+b\cdot Du=0$.
\end{theorem}
\begin{proof}
\leanok
\sketch Write $u=g\circ L$ for $L(x,t)=x-tb$. The transport operator is the
directional derivative along $(b,1)$, and $L(b,1)=0$; the chain rule therefore
gives $u_t+b\cdot Du=0$.
\end{proof}

\begin{theorem}[Initial condition for the transport solution]
\label{thm:transport-initial-condition}
\lean{EvansLib.transportSolution_init}
\uses{def:transport-explicit-solution}
\leanok
\dcref{ch2:2.1.1-thm-2}
\group{Transport Equation}



The explicit solution satisfies $u(x,0)=g(x)$ for every $x\in\R^n$.
\end{theorem}
\begin{proof}
\leanok
\sketch Substitute $t=0$ in the defining formula.
\end{proof}

\begin{theorem}[Uniqueness for the transport initial-value problem]
\label{thm:transport-ivp-uniqueness}
\lean{EvansLib.transportSolution_unique}
\uses{def:transport-explicit-solution}
\leanok
\dcref{ch2:2.1.1-thm-3}
\group{Transport Equation}



If a differentiable $v$ satisfies $v_t+b\cdot Dv=0$ on $\R^n\times\R$ and
$v(x,0)=g(x)$, then $v(x,t)=g(x-tb)$.
\end{theorem}
\begin{proof}
\leanok
\sketch For fixed $(x,t)$, the function
$z(s)=v(x+sb,t+s)$ has derivative $z'(s)=v_t+b\cdot Dv=0$. Comparing $s=0$
with the intersection $s=-t$ of the characteristic and the initial plane gives
$v(x,t)=v(x-tb,0)=g(x-tb)$.
\end{proof}

\begin{remark}[Weak solutions of the transport equation]
\label{rem:transport-equation-weak-solution-discontinuous-data}
\dcref{ch2:2.1.1-rem-1}
\group{Transport Equation}
For nonsmooth, even discontinuous, initial data $g$, the same formula
$u(x,t)=g(x-tb)$ defines the natural weak transport solution, although it need
not be $C^1$.
\end{remark}

\subsection{Nonhomogeneous Problem}
\label{sec:transport-equation-nonhomogeneous-problem}

For
$$(4)\qquad
\begin{cases}
u_t+b\cdot Du=f&\text{in }\R^n\times(0,\infty),\\
u=g&\text{on }\R^n\times\{t=0\},
\end{cases}$$
integration along characteristics gives the following formula.

\begin{definition}[Explicit solution of the nonhomogeneous transport equation]
\label{def:transport-nonhom-solution}
\lean{EvansLib.transportSolutionNonhom}
\uses{def:transport-explicit-solution}
\leanok
\dcref{ch2:2.1.2-def-1}
\group{Transport Equation}



For $b\in\R^n$, $g:\R^n\to\R$, and $f:\R^n\times\R\to\R$, define
$$u(x,t):=g(x-tb)+\int_0^t f\bigl(x+(s-t)b,s\bigr)\,ds.$$
The integrand samples $f$ on the characteristic through $(x,t)$ at time $s$.
\end{definition}

\begin{theorem}[The nonhomogeneous solution attains the initial data]
\label{thm:transport-nonhom-initial-condition}
\lean{EvansLib.transportSolutionNonhom_init}
\uses{def:transport-nonhom-solution}
\leanok
\dcref{ch2:2.1.2-thm-1}
\group{Transport Equation}



The nonhomogeneous solution satisfies $u(x,0)=g(x)$.
\end{theorem}
\begin{proof}
\leanok
\sketch At $t=0$ the integral vanishes and $g(x-tb)=g(x)$.
\end{proof}

\begin{theorem}[The nonhomogeneous transport equation holds along characteristics]
\label{thm:transport-nonhom-solves}
\lean{EvansLib.transportSolutionNonhom_char_deriv}
\uses{def:transport-nonhom-solution}
\leanok
\dcref{ch2:2.1.2-thm-2}
\group{Transport Equation}



If $f$ is continuous, then
$$\left.\frac{d}{dr}\right|_{r=0}u\bigl((x,t)+r(b,1)\bigr)=f(x,t).$$
If $u$ is differentiable, this is $u_t+b\cdot Du=f$.
\end{theorem}
\begin{proof}
\leanok
\sketch Along $(b,1)$ the homogeneous term is constant, while the integral
becomes an integral with fixed integrand and upper endpoint $t+r$. The
fundamental theorem of calculus gives the displayed derivative.
\end{proof}

\begin{remark}[Method of characteristics]
\label{rem:transport-equation-method-of-characteristics}
\dcref{ch2:2.1.2-rem-1}
\group{Transport Equation}
The preceding formulas reduce the transport PDE to an ODE along each
characteristic line.
\end{remark}

\section{Laplace's Equation}
\label{sec:laplaces-equation}

Let $U\subset\R^n$ be open. Laplace's and Poisson's equations are
$$(1)\qquad \Delta u=0,
\qquad
(2)\qquad -\Delta u=f,$$
where $\Delta u=\sum_{i=1}^n u_{x_i x_i}$.

\begin{definition}[Harmonic function]
\label{def:harmonic-function}
\lean{InnerProductSpace.HarmonicOnNhd}
\dcref{ch2:2.2-def-1}
\group{Fundamental Solutions}

\mathlibok
A function $u\in C^2(U)$ is \emph{harmonic} on $U$ if $\Delta u=0$ on $U$.
\end{definition}

\subsection{Fundamental Solution}
\label{sec:laplace-fundamental-solution}

\subsubsection{Derivation of fundamental solution}
\label{sec:laplace-fundamental-solution-derivation}

For a radial function $u(x)=v(r)$, $r=|x|>0$,
$$\Delta u=v''(r)+\frac{n-1}{r}v'(r).$$
Thus radial harmonic functions have the form $b\log r+c$ for $n=2$ and
$br^{2-n}+c$ for $n\ge3$.

\begin{definition}[Fundamental solution of Laplace's equation]
\label{def:fundamental-solution-laplace}
\lean{EvansLib.laplaceFund}
\uses{def:harmonic-function}
\leanok
\dcref{ch2:2.2.1-def-1}
\group{Fundamental Solutions}



For $x\ne0$, define
$$(6)\qquad
\Phi(x):=
\begin{cases}
-\dfrac{1}{2\pi}\log|x|&(n=2),\\[3pt]
\dfrac{1}{n(n-2)\alpha(n)}|x|^{2-n}&(n\ge3),
\end{cases}$$
where $\alpha(n)$ is the volume of the unit ball in $\R^n$.
\end{definition}

\begin{lemma}[The fundamental solution is harmonic away from the origin]
\label{lem:laplace-fundamental-solution-harmonic}
\lean{EvansLib.laplaceFund_harmonic,EvansLib.laplaceFund_harmonicOnNhd}
\uses{def:fundamental-solution-laplace, def:harmonic-function}
\leanok
\dcref{ch2:2.2.1-lem-1}
\group{Fundamental Solutions}



For every $x\ne0$, $\Delta\Phi(x)=0$.
\end{lemma}
\begin{proof}
\leanok
\sketch Write $\Phi(x)=g(|x|^2)$. The chain rule gives
$$\Delta\Phi(x)=4|x|^2g''(|x|^2)+2n g'(|x|^2).$$
Substitution of $g(\rho)=-(4\pi)^{-1}\log\rho$ for $n=2$ or
$g(\rho)=[n(n-2)\alpha(n)]^{-1}\rho^{(2-n)/2}$ for $n\ge3$ makes this
expression vanish.
\end{proof}

The radial formula also yields, for a constant $C>0$,
$$(7)\qquad
|D\Phi(x)|\le \frac{C}{|x|^{n-1}},
\qquad
|D^2\Phi(x)|\le \frac{C}{|x|^n}
\quad(x\ne0).$$

\subsubsection{Poisson's equation}
\label{sec:laplace-poissons-equation}

For $f\in C_c^2(\R^n)$, set
$$(8)\qquad u(x):=(\Phi*f)(x)=\int_{\R^n}\Phi(x-y)f(y)\,dy.$$

\begin{theorem}[Solving Poisson's equation]
\label{thm:poisson-equation-solution-formula}
\uses{def:fundamental-solution-laplace}
\dcref{ch2:2.2-thm-1}
\group{Fundamental Solutions}

The function in (8) belongs to $C^2(\R^n)$ and satisfies
$$-\Delta u=f\quad\text{in }\R^n.$$
\end{theorem}
\begin{proof}
\sketch Rewrite $u(x)=\int\Phi(y)f(x-y)\,dy$; local integrability of $\Phi$
and compact support of $f$ justify differentiating $f$ under the integral and
give $u\in C^2$. Split $\Delta u=I_\varepsilon+J_\varepsilon$ over
$B(0,\varepsilon)$ and its complement. The estimates
$$|I_\varepsilon|\le
\begin{cases}C\varepsilon^2|\log\varepsilon|&n=2,\\
C\varepsilon^2&n\ge3,
\end{cases}$$
and, after one integration by parts, a boundary error bounded by
$C\varepsilon|\log\varepsilon|$ for $n=2$ and $C\varepsilon$ for $n\ge3$,
make these terms vanish. A second integration by parts uses
$\Delta\Phi=0$ away from $0$ and
$$\frac{\partial\Phi}{\partial\nu}
=\frac{1}{n\alpha(n)\varepsilon^{n-1}}
\quad\text{on }\partial B(0,\varepsilon),$$
where $\nu$ points into the deleted ball. The remaining boundary average tends
to $-f(x)$, hence $\Delta u=-f$.
\end{proof}

\begin{remark}[Dirac measure notation]
\label{rem:poisson-equation-dirac-delta-notation}
\uses{thm:poisson-equation-solution-formula}
\dcref{ch2:2.2.1-rem-1}
\group{Fundamental Solutions}

Distributionally,
$$-\Delta\Phi=\delta_0,$$
so $-\Delta(\Phi*f)=f$. The conclusion of
\cref{thm:poisson-equation-solution-formula} remains valid under weaker
hypotheses on $f$; see \cite{GilbargTrudinger1983}.
\end{remark}

\subsection{Mean-Value Formulas}
\label{sec:laplace-mean-value-formulas}

\begin{theorem}[Mean-value formulas for Laplace's equation]
\label{thm:mean-value-formulas-laplace}
\lean{EvansLib.unitSphereRadialIntegralAt_eq_volume_mul_of_harmonicOnNhd, EvansLib.HarmonicOnNhd.hasBallMeanValueProperty}
\uses{def:harmonic-function}
\dcref{ch2:2.2-thm-2}
\group{Harmonic Functions}

If $u\in C^2(U)$ is harmonic, then for every ball $B(x,r)\subset U$,
$$(16)\qquad
u(x)=\fint_{\partial B(x,r)}u\,dS
=\fint_{B(x,r)}u\,dy.$$
\end{theorem}
\begin{proof}
\sketch Let $\phi(r)=\fint_{\partial B(x,r)}u\,dS$. Differentiation and the
divergence theorem give
$$\phi'(r)=\frac{r}{n}\fint_{B(x,r)}\Delta u\,dy=0.$$
Thus $\phi(r)=\lim_{s\downarrow0}\phi(s)=u(x)$. Integrating the spherical
identity in polar coordinates yields the ball identity.
\end{proof}

\begin{theorem}[Converse to mean-value property]
\label{thm:converse-mean-value-property-laplace}
\lean{EvansLib.harmonicOnNhd_of_unitSphereRadialIntegralAt_eq_zero}
\uses{def:harmonic-function}
\leanok
\dcref{ch2:2.2-thm-3}
\group{Harmonic Functions}

If $u\in C^2(U)$ and
$$u(x)=\fint_{\partial B(x,r)}u\,dS$$
for every $B(x,r)\subset U$, then $u$ is harmonic.
\end{theorem}
\begin{proof}
\sketch The preceding differentiation formula makes every spherical average
constant in $r$. If $\Delta u(x_0)\ne0$, after replacing $u$ by $-u$ if needed,
$\Delta u>0$ on a small ball about $x_0$, contradicting
$0=\phi'(r)=\frac r n\fint_{B(x_0,r)}\Delta u$.
\end{proof}

\begin{theorem}[Converse to mean-value property, ball form]
\label{thm:converse-mean-value-property-mvp}
\lean{EvansLib.HasBallMeanValueProperty.laplacian_eq_zero, EvansLib.HasBallMeanValueProperty.harmonicOnNhd}
\uses{def:ball-mean-value-property, def:harmonic-function}
\leanok
\group{Harmonic Functions}



Let $U\subset\R^n$ be open. For continuous functions on $U$, harmonicity is
equivalent to the ball mean-value property, and every such function is smooth.
\end{theorem}
\begin{proof}
\uses{thm:harmonic-functions-smooth}
\leanok
\sketch By \cref{thm:harmonic-functions-smooth}, $u$ is smooth. At $x\in U$,
write $L=Du(x)$ and $B=D^2u(x)$. Averaging the second-order Taylor remainder
over $B(0,1)$ and using the mean-value property and reflection symmetry gives
$$\frac1{r^2}\fint_{B(0,1)}
\left[u(x+rw)-u(x)-rL(w)-\frac{r^2}{2}B(w,w)\right]dw
=-\frac12\fint_{B(0,1)}B(w,w)\,dw.$$
Dominated convergence sends the left side to $0$. Off-diagonal moments vanish
and all diagonal moments agree, so the right side is a nonzero constant times
$\operatorname{tr}B=\Delta u(x)$. Thus $\Delta u(x)=0$.
\end{proof}

\subsubsection{Consequences of the mean-value property}
\label{sec:laplace-mean-value-consequences}

\begin{definition}[Ball mean-value property]
\label{def:ball-mean-value-property}
\lean{EvansLib.HasBallMeanValueProperty}
\leanok
\group{Harmonic Functions}


A function $u:\R^n\to\R$ has the \emph{ball mean-value property} on an open
set $U$ if
$$u(x)=\fint_{B(x,r)}u\,dy$$
whenever $r>0$ and $\overline{B(x,r)}\subset U$.
\end{definition}

\begin{theorem}[Strong maximum principle, mean-value form]
\label{thm:strong-maximum-principle-mvp}
\lean{EvansLib.eqOn_of_isPreconnected_of_isMaxOn, EvansLib.exists_frontier_isMaxOn}
\uses{def:ball-mean-value-property}
\leanok
\group{Maximum Principles}



Let $n\ge1$ and let $u$ have the ball mean-value property on an open set
$U\subset\R^n$.

(i) If $U$ is bounded and nonempty and $u\in C(\overline U)$, then $u$ attains
its maximum on $\partial U$.

(ii) If $U$ is connected, $u\in C(U)$, and $u$ attains its maximum at a point
of $U$, then $u$ is constant.
\end{theorem}
\begin{proof}
\leanok
\sketch If $u(y)\le u(x)$ on $B(x,r)$, the mean-value identity makes the
continuous nonnegative function $u(x)-u$ have zero integral, hence vanish on
the ball. For (ii), the set of interior maximizers is nonempty, relatively
closed, and open by this local fact; connectedness makes it all of $U$. For
(i), compactness gives a maximizer, and the same local fact prevents all
maximizers from forming a nonempty bounded subset both open and closed in
$\R^n$ unless one lies on $\partial U$.
\end{proof}

\begin{theorem}[Uniqueness for the Dirichlet problem, mean-value form]
\label{thm:uniqueness-dirichlet-mvp}
\lean{EvansLib.eqOn_closure_of_eqOn_frontier}
\uses{def:ball-mean-value-property, thm:strong-maximum-principle-mvp}
\group{Maximum Principles}



Let $U\subset\R^n$ be bounded and open. If $u,v\in C(\overline U)$ have the
ball mean-value property on $U$ and agree on $\partial U$, then
$u=v$ on $\overline U$.
\end{theorem}
\begin{proof}
\uses{thm:strong-maximum-principle-mvp}
\leanok
\sketch The functions $w=u-v$ and $-w$ have the ball mean-value property and
vanish on the boundary. Applying the maximum principle to both gives
$w\le0$ and $w\ge0$.
\end{proof}

\begin{theorem}[Harnack's inequality, local form]
\label{thm:harnack-local-mvp}
\lean{EvansLib.harnack_local}
\uses{def:ball-mean-value-property}
\group{Harmonic Functions}



Let $n\ge1$, let $u\in C(U)$ be nonnegative with the ball mean-value property,
and suppose $|x-y|\le r$ and $\overline{B(x,2r)}\subset U$. Then
$$u(y)\le2^n u(x).$$
\end{theorem}
\begin{proof}
\leanok
\sketch Since $B(y,r)\subset B(x,2r)$ and $u\ge0$,
$$u(x)=\fint_{B(x,2r)}u
\ge\frac{|B(y,r)|}{|B(x,2r)|}\fint_{B(y,r)}u
=2^{-n}u(y).$$
\end{proof}

\begin{theorem}[Harnack's inequality, mean-value form]
\label{thm:harnack-global-mvp}
\lean{EvansLib.exists_harnack_const}
\uses{def:ball-mean-value-property, thm:harnack-local-mvp}
\leanok
\group{Harmonic Functions}



Let $n\ge1$, let $U\subset\R^n$ be open, and let $V$ be connected and bounded
with $\overline V\subset U$. There is $C>0$, depending only on $V$, such that
$$u(x)\le C u(y)\qquad(x,y\in V)$$
for every nonnegative $u\in C(U)$ with the ball mean-value property.
\end{theorem}
\begin{proof}
\uses{thm:harnack-local-mvp}
\leanok
\sketch Choose $r>0$ so every $B(z,2r)$, $z\in\overline V$, lies in $U$, and
cover $\overline V$ by finitely many $r/2$-balls. Connectivity makes the
intersection graph of those meeting $V$ connected. Iterating the local bound
along a path of at most $N$ balls yields $u(x)\le(2^n)^N u(y)$.
\end{proof}

\begin{theorem}[Liouville's theorem, mean-value form]
\label{thm:liouville-mvp}
\lean{EvansLib.liouville_of_bddBelow, EvansLib.liouville_of_isBounded}
\uses{def:ball-mean-value-property, thm:harnack-local-mvp}
\leanok
\group{Harmonic Functions}



Let $n\ge1$. If $u\in C(\R^n)$ has the ball mean-value property and is bounded
below, then $u$ is constant. In particular, every bounded such function is
constant.
\end{theorem}
\begin{proof}
\uses{thm:harnack-local-mvp}
\leanok
\sketch Put $m=\inf_{\R^n}u$ and $v=u-m\ge0$. The local Harnack estimate is
available at every scale, so $v(x)\le2^n v(y)$ for all $x,y$. Taking the
infimum over $y$ forces $v(x)=0$.
\end{proof}

\subsection{Properties of Harmonic Functions}
\label{sec:laplace-properties-harmonic-functions}

\subsubsection{Strong maximum principle, uniqueness}
\label{sec:laplace-strong-maximum-principle}

\begin{theorem}[Strong maximum principle]
\label{thm:strong-maximum-principle-laplace}
\lean{EvansLib.harmonic_exists_frontier_isMaxOn, EvansLib.harmonic_eqOn_of_isPreconnected_of_isMaxOn}
\uses{thm:mean-value-formulas-laplace}
\leanok
\dcref{ch2:2.2-thm-4}
\group{Maximum Principles}

Let $U$ be bounded and open, and let
$u\in C^2(U)\cap C(\overline U)$ be harmonic.

(i) $\max_{\overline U}u=\max_{\partial U}u$.

(ii) If $U$ is connected and $u$ attains its maximum at an interior point,
then $u$ is constant on $U$.
\end{theorem}
\begin{proof}
\sketch At an interior maximizer $x_0$, the mean-value formula forces
$u\equiv u(x_0)$ on every ball centered at $x_0$ and compactly contained in
$U$. The set of maximizers is therefore open and closed in $U$, proving (ii);
(i) follows unless $u$ is constant. Apply the result to $-u$ for minima.
\end{proof}

\begin{remark}[Strict positivity from boundary data]
\label{rem:strong-maximum-principle-positivity}
\lean{EvansLib.harmonic_pos_of_nonneg_on_frontier}
\uses{thm:strong-maximum-principle-laplace}
\leanok
\dcref{ch2:2.2.3-rem-2}
\group{Maximum Principles}

If $U$ is connected, $u\in C^2(U)\cap C(\overline U)$ is harmonic, and its
boundary data are nonnegative and positive somewhere, then $u>0$ throughout
$U$.
\end{remark}

\begin{theorem}[Uniqueness]
\label{thm:uniqueness-poisson-dirichlet}
\lean{EvansLib.eqOn_closure_of_eqOn_frontier_of_laplacian_eq}
\uses{thm:strong-maximum-principle-laplace}
\leanok
\dcref{ch2:2.2-thm-5}
\group{Maximum Principles}

For $g\in C(\partial U)$ and $f\in C(U)$, there is at most one
$u\in C^2(U)\cap C(\overline U)$ satisfying
$$(17)\qquad
\begin{cases}
-\Delta u=f&\text{in }U,\\
u=g&\text{on }\partial U.
\end{cases}$$
\end{theorem}
\begin{proof}
\sketch The difference of two solutions is harmonic and vanishes on
$\partial U$. Apply the maximum principle to the difference and its negative.
\end{proof}

\subsubsection{Regularity}
\label{sec:laplace-regularity}

\begin{theorem}[Smoothness]
\label{thm:harmonic-functions-smooth}
\lean{EvansLib.HasBallMeanValueProperty.contDiffOn}
\uses{def:ball-mean-value-property, thm:mean-value-formulas-laplace}
\leanok
\dcref{ch2:2.2-thm-6}
\group{Harmonic Functions}



If $u\in C(U)$ has the mean-value property (16) on every ball compactly
contained in $U$, then $u\in C^\infty(U)$.
\end{theorem}
\begin{proof}
\leanok
\sketch Let $\eta$ be a radial standard mollifier and set
$u^\varepsilon=\eta_\varepsilon*u$ on
$U_\varepsilon=\{x\in U:\dist(x,\partial U)>\varepsilon\}$. Polar
coordinates and the spherical mean-value formula give
$$u^\varepsilon(x)
=\frac1{\varepsilon^n}\int_0^\varepsilon
\eta(r/\varepsilon)\left(\int_{\partial B(x,r)}u\,dS\right)dr
=u(x)\int\eta_\varepsilon=u(x).$$
Thus $u$ agrees locally with a smooth mollification.
\end{proof}

\subsubsection{Local estimates for derivatives}
\label{sec:laplace-local-estimates-derivatives}

\begin{theorem}[Estimates on derivatives]
\label{thm:harmonic-function-derivative-estimates}
\uses{thm:mean-value-formulas-laplace}
\dcref{ch2:2.2-thm-7}
\group{Harmonic Functions}

If $u$ is harmonic in $U$, $B(x_0,r)\subset U$, and $|\alpha|=k$, then
$$(18)\qquad
|D^\alpha u(x_0)|\le
\frac{C_k}{r^{n+k}}\|u\|_{L^1(B(x_0,r))},$$
where
$$(19)\qquad
C_0=\frac1{\alpha(n)},
\qquad
C_k=\frac{(2^{n+1}nk)^k}{\alpha(n)}\quad(k\ge1).$$
\end{theorem}
\begin{proof}
\sketch Induct on $k$. The mean-value formula gives $k=0$. Since derivatives
of a harmonic function are harmonic, averaging $u_{x_i}$ over $B(x_0,r/2)$
and integrating by parts yields
$$|u_{x_i}(x_0)|\le\frac{2n}{r}
\|u\|_{L^\infty(\partial B(x_0,r/2))},$$
and a radius-$r/2$ mean-value estimate gives the asserted constant for $k=1$.
For $|\alpha|=k$, write $D^\alpha u=(D^\beta u)_{x_i}$ with
$|\beta|=k-1$, average on the sphere of radius $(k-1)r/k$, and apply the
induction hypothesis on radius-$r/k$ balls. Combining the factors gives
$C_k=(2^{n+1}nk)^k/\alpha(n)$.
\end{proof}

\subsubsection{Liouville's Theorem}
\label{sec:laplace-liouvilles-theorem}

\begin{theorem}[Liouville's Theorem]
\label{thm:liouville-theorem-harmonic}
\lean{EvansLib.harmonic_liouville_of_isBounded}
\uses{thm:harmonic-function-derivative-estimates, thm:liouville-mvp}
\leanok
\dcref{ch2:2.2-thm-8}
\group{Harmonic Functions}

Every bounded harmonic function $u:\R^n\to\R$ is constant.
\end{theorem}
\begin{proof}
\sketch For any $x_0$ and $r>0$, the derivative estimate gives
$$|Du(x_0)|\le\frac{C_1\alpha(n)}r\|u\|_{L^\infty(\R^n)}.$$
Letting $r\to\infty$ gives $Du=0$.
\end{proof}

\begin{theorem}[Representation formula]
\label{thm:representation-formula-bounded-solutions}
\uses{def:fundamental-solution-laplace, thm:liouville-theorem-harmonic}
\dcref{ch2:2.2-thm-9}
\group{Harmonic Functions}

Let $n\ge3$ and $f\in C_c^2(\R^n)$. Every bounded solution of
$-\Delta u=f$ on $\R^n$ has the form
$$u(x)=\int_{\R^n}\Phi(x-y)f(y)\,dy+C$$
for a constant $C$.
\end{theorem}
\begin{proof}
\sketch Since $\Phi(x)\to0$ as $|x|\to\infty$, the convolution is a bounded
solution. Its difference from any other bounded solution is bounded and
harmonic, hence constant by Liouville's theorem.
\end{proof}

\begin{remark}
\label{rem:representation-formula-n2-unbounded}
\uses{def:fundamental-solution-laplace, thm:representation-formula-bounded-solutions}
\dcref{ch2:2.2.3-rem-3}
\group{Harmonic Functions}

For $n=2$, $\Phi(x)=-(2\pi)^{-1}\log|x|$ is unbounded at infinity, and so may
be $\Phi*f$.
\end{remark}

\subsubsection{Analyticity}
\label{sec:laplace-analyticity}

\begin{theorem}[Analyticity]
\label{thm:harmonic-functions-analytic}
\uses{thm:harmonic-function-derivative-estimates, thm:harmonic-functions-smooth}
\dcref{ch2:2.2-thm-10}
\group{Harmonic Functions}

Every harmonic function is analytic in its domain.
\end{theorem}
\begin{proof}
\sketch Fix $x_0\in U$, put $r=\dist(x_0,\partial U)/4$, and set
$M=[\alpha(n)r^n]^{-1}\|u\|_{L^1(B(x_0,2r))}$. The derivative estimate,
Stirling's formula [RD, \S8.22], and the multinomial theorem give
$$(22)\qquad
\|D^\alpha u\|_{L^\infty(B(x_0,r))}
\le CM\left(\frac{2^{n+1}n^2e}{r}\right)^{|\alpha|}\alpha!.$$
Taylor's formula along $t\mapsto u(x_0+t(x-x_0))$ then bounds the degree-$N$
remainder by $CM2^{-N}$ whenever
$|x-x_0|<r/(2^{n+2}n^3e)$. Hence the Taylor series converges to $u$ near
$x_0$.
\end{proof}

\subsubsection{Harnack's inequality}
\label{sec:laplace-harnacks-inequality}

\begin{theorem}[Harnack's inequality]
\label{thm:harnacks-inequality}
\lean{EvansLib.exists_harnack_const_harmonic}
\uses{thm:mean-value-formulas-laplace, thm:harnack-global-mvp}
\leanok
\dcref{ch2:2.2-thm-11}
\group{Harmonic Functions}

For every connected open $V\Subset U$, there is $C>0$, depending only on
$V$, such that
$$\sup_Vu\le C\inf_Vu$$
for every nonnegative harmonic function $u$ in $U$.
\end{theorem}
\begin{proof}
\sketch With $r=\dist(V,\partial U)/4$, the mean-value formula and
$B(y,r)\subset B(x,2r)$ show that $2^{-n}u(y)\le u(x)\le2^n u(y)$ whenever
$x,y\in V$ and $|x-y|\le r$. A finite chain of overlapping balls covering the
compact connected set $\overline V$ yields the global constant.
\end{proof}

\subsection{Green's Function}
\label{sec:laplace-greens-function}

Let $U\subset\R^n$ be bounded with $C^1$ boundary. For $x\in U$, let
$\varphi^x$ solve
$$\Delta\varphi^x=0\ \text{in }U,
\qquad
\varphi^x(y)=\Phi(y-x)\ \text{on }\partial U.$$

\subsubsection{Derivation of Green's function}
\label{sec:laplace-greens-function-derivation}

\begin{definition}[Green's function]
\label{def:greens-function}
\uses{def:fundamental-solution-laplace}
\dcref{ch2:2.2.4-def-1}
\group{Green's Functions}

Green's function for $U$ is
$$G(x,y):=\Phi(y-x)-\varphi^x(y)
\qquad(x,y\in U,\ x\ne y).$$
\end{definition}

Green's identity on $U\setminus B(x,\varepsilon)$, followed by
$\varepsilon\downarrow0$, gives for $u\in C^2(\overline U)$
$$u(x)=-\int_{\partial U}u(y)\frac{\partial G}{\partial\nu}(x,y)\,dS(y)
-\int_U G(x,y)\Delta u(y)\,dy.$$

\begin{theorem}[Representation formula using Green's function]
\label{thm:representation-formula-greens-function}
\uses{def:greens-function}
\dcref{ch2:2.2-thm-12}
\group{Green's Functions}

If $u\in C^2(\overline U)$ satisfies
$$-\Delta u=f\ \text{in }U,
\qquad
u=g\ \text{on }\partial U,$$
then
$$(30)\qquad
u(x)=-\int_{\partial U}g(y)\frac{\partial G}{\partial\nu}(x,y)\,dS(y)
+\int_U f(y)G(x,y)\,dy.$$
\end{theorem}

\begin{remark}[Green's function and the Dirac measure]
\label{rem:greens-function-dirac-delta}
\uses{def:greens-function}
\dcref{ch2:2.2.4-rem-3}
\group{Green's Functions}

For fixed $x\in U$, Green's function satisfies symbolically
$$-\Delta_yG(x,\cdot)=\delta_x\ \text{in }U,
\qquad
G(x,\cdot)=0\ \text{on }\partial U.$$
\end{remark}

\begin{theorem}[Symmetry of Green's function]
\label{thm:greens-function-symmetry}
\uses{def:greens-function}
\dcref{ch2:2.2-thm-13}
\group{Green's Functions}

For distinct $x,y\in U$,
$$G(y,x)=G(x,y).$$
\end{theorem}
\begin{proof}
\sketch Apply Green's identity to $G(x,\cdot)$ and $G(y,\cdot)$ on
$U\setminus(B(x,\varepsilon)\cup B(y,\varepsilon))$. The outer boundary terms
vanish. As $\varepsilon\downarrow0$, the two small-sphere integrals converge
to $G(y,x)$ and $G(x,y)$, respectively.
\end{proof}

\subsubsection{Green's function for a half-space}
\label{sec:laplace-greens-function-half-space}

Write $\R^n_+=\{x\in\R^n:x_n>0\}$.

\begin{definition}[Reflection in a half-space]
\label{def:reflection-half-space}
\lean{EvansLib.reflectHalfSpace}
\leanok
\dcref{ch2:2.2.4-def-2}
\group{Green's Functions}


For $x=(x_1,\ldots,x_{n-1},x_n)\in\R^n_+$, define
$$\overline x=(x_1,\ldots,x_{n-1},-x_n).$$
\end{definition}

\begin{definition}[Green's function for the half-space]
\label{def:greens-function-half-space}
\lean{EvansLib.greensHalfSpace}
\uses{def:greens-function, def:fundamental-solution-laplace, def:reflection-half-space}
\dcref{ch2:2.2.4-def-3}
\group{Green's Functions}



Green's function for $\R^n_+$ is
$$G(x,y):=\Phi(y-x)-\Phi(y-\overline x)
\qquad(x,y\in\R^n_+,\ x\ne y).$$
It vanishes when $y\in\partial\R^n_+$.
\end{definition}

On $y\in\partial\R^n_+$,
$$\frac{\partial G}{\partial\nu}(x,y)
=-\frac{2x_n}{n\alpha(n)|x-y|^n},$$
so the Poisson kernel is
$$K(x,y)=\frac{2x_n}{n\alpha(n)|x-y|^n}.$$

\begin{theorem}[Poisson's formula for half-space]
\label{thm:poissons-formula-half-space}
\uses{def:greens-function-half-space}
\dcref{ch2:2.2-thm-14}
\group{Green's Functions}

Let $g\in C(\R^{n-1})\cap L^\infty(\R^{n-1})$ and define
$$u(x)=\frac{2x_n}{n\alpha(n)}
\int_{\partial\R^n_+}\frac{g(y)}{|x-y|^n}\,dy.$$
Then $u\in C^\infty(\R^n_+)\cap L^\infty(\R^n_+)$,
$\Delta u=0$ in $\R^n_+$, and
$$\lim_{\substack{x\to x^0\\x\in\R^n_+}}u(x)=g(x^0)
\qquad(x^0\in\partial\R^n_+).$$
\end{theorem}
\begin{proof}
\sketch The kernel is harmonic in $x$ and has boundary integral $1$, giving
boundedness, smoothness, and harmonicity under the integral. For the boundary
limit, split $u(x)-g(x^0)$ into $|y-x^0|<\delta$ and its complement.
Continuity controls the first integral; on the second,
$|x-y|\ge|y-x^0|/2$, so its absolute value is bounded by a constant times
$$x_n\int_{|y-x^0|\ge\delta}|y-x^0|^{-n}\,dy\to0.$$
\end{proof}

\subsubsection{Green's function for a ball}
\label{sec:laplace-greens-function-ball}

\begin{definition}[Dual point]
\label{def:dual-point-unit-sphere}
\lean{EvansLib.dualPoint}
\leanok
\dcref{ch2:2.2.4-def-4}
\group{Green's Functions}


For $x\ne0$, the point
$$\overline x:=\frac{x}{|x|^2}$$
is dual to $x$ under inversion through $\partial B(0,1)$.
\end{definition}

For $|y|=1$, one has $|x|^2|y-\overline x|^2=|x-y|^2$.

\begin{definition}[Green's function for the unit ball]
\label{def:greens-function-ball}
\lean{EvansLib.greensBall}
\uses{def:greens-function, def:fundamental-solution-laplace, def:dual-point-unit-sphere}
\dcref{ch2:2.2.4-def-5}
\group{Green's Functions}



Green's function for the unit ball is
$$(41)\qquad
G(x,y):=\Phi(y-x)-\Phi\bigl(|x|(y-\overline x)\bigr)
\qquad(x,y\in B(0,1),\ x\ne y).$$
This formula is valid for $n=2$ and $n\ge3$, and $G(x,y)=0$ for
$y\in\partial B(0,1)$.
\end{definition}

For $B(0,r)$ the Poisson kernel is
$$K(x,y)=\frac{r^2-|x|^2}{n\alpha(n)r|x-y|^n},$$
and Poisson's formula is
$$(45)\qquad
u(x)=\frac{r^2-|x|^2}{n\alpha(n)r}
\int_{\partial B(0,r)}\frac{g(y)}{|x-y|^n}\,dS(y).$$

\begin{theorem}[Poisson's formula for ball]
\label{thm:poissons-formula-ball}
\uses{def:greens-function-ball}
\dcref{ch2:2.2-thm-15}
\group{Green's Functions}

If $g\in C(\partial B(0,r))$ and $u$ is defined by (45), then
$u\in C^\infty(B^0(0,r))$, $\Delta u=0$ there, and
$$\lim_{\substack{x\to x^0\\x\in B^0(0,r)}}u(x)=g(x^0)
\qquad(x^0\in\partial B(0,r)).$$
\end{theorem}

\subsection{Energy Methods}
\label{sec:laplace-energy-methods}

\subsubsection{Uniqueness}
\label{sec:laplace-energy-uniqueness}

Let $U$ be bounded with $C^1$ boundary and consider
$$(46)\qquad
\begin{cases}
-\Delta u=f&\text{in }U,\\
u=g&\text{on }\partial U.
\end{cases}$$

\begin{theorem}[Uniqueness]
\label{thm:uniqueness-poisson-energy-method}
\lean{EvansLib.eqOn_closure_of_eqOn_frontier_of_laplacian_eq}
\leanok
\dcref{ch2:2.2-thm-16}
\group{Energy Methods}
There is at most one $u\in C^2(\overline U)$ satisfying (46).
\end{theorem}
\begin{proof}
  \uses{thm:uniqueness-poisson-dirichlet}
\sketch The difference $w$ of two solutions is harmonic and vanishes on the
boundary. Integration by parts gives
$$0=-\int_Uw\Delta w\,dx=\int_U|Dw|^2\,dx,$$
so $w$ is constant and hence zero.
\end{proof}

\subsubsection{Dirichlet's principle}
\label{sec:laplace-dirichlets-principle}

Define
$$I[w]=\int_U\left(\frac12|Dw|^2-wf\right)dx,
\qquad
\mathcal A=\{w\in C^2(\overline U):w=g\text{ on }\partial U\}.$$

\begin{theorem}[Dirichlet's principle]
\label{thm:dirichlets-principle}
\dcref{ch2:2.2-thm-17}
\group{Energy Methods}
A function $u\in\mathcal A$ solves (46) if and only if
$$(47)\qquad I[u]=\min_{w\in\mathcal A}I[w].$$
\end{theorem}
\begin{proof}
\sketch If $u$ solves (46), test against $u-w$ and integrate by parts. The
identity
$$\int_U|Du|^2-uf=\int_UDu\cdot Dw-wf$$
and $Du\cdot Dw\le(|Du|^2+|Dw|^2)/2$ give $I[u]\le I[w]$. Conversely, for
$v\in C_c^\infty(U)$, minimality makes the derivative at $\tau=0$ of
$I[u+\tau v]$ vanish:
$$0=\int_U Du\cdot Dv-vf=\int_U(-\Delta u-f)v.$$
Thus $-\Delta u=f$.
\end{proof}

\section{Heat Equation}
\label{sec:heat-equation}

Let $U\subset\R^n$ be open, $t>0$, and $D=D_x$. The homogeneous and
nonhomogeneous heat equations are
$$(1)\qquad u_t-\Delta u=0,
\qquad
(2)\qquad u_t-\Delta u=f.$$

\subsection{Fundamental Solution}
\label{sec:heat-fundamental-solution}

\subsubsection{Derivation of fundamental solution}
\label{sec:heat-fundamental-solution-derivation}

The heat operator is invariant under $(x,t)\mapsto(\lambda x,\lambda^2t)$.
A self-similar radial ansatz
$u(x,t)=t^{-n/2}w(|x|t^{-1/2})$ reduces the equation to
$w'=-(r/2)w$, hence $w(r)=be^{-r^2/4}$.

\begin{definition}[Fundamental solution of the heat equation]
\label{def:fundamental-solution-heat-equation}
\lean{EvansLib.heatKernel}
\leanok
\dcref{ch2:2.3.1-def-1}
\group{Heat Equation}


The fundamental solution of the heat equation is
$$\Phi(x,t):=
\begin{cases}
(4\pi t)^{-n/2}e^{-|x|^2/(4t)}&t>0,\\
0&t<0.
\end{cases}$$
\end{definition}

\begin{lemma}[Integral of fundamental solution]
\label{lem:integral-fundamental-solution-heat}
\lean{EvansLib.heatKernelSpatial_integral}
\uses{def:fundamental-solution-heat-equation}
\leanok
\dcref{ch2:2.3.1-lem-1}
\group{Heat Equation}



For every $t>0$,
$$\int_{\R^n}\Phi(x,t)\,dx=1.$$
\end{lemma}
\begin{proof}
\leanok
\sketch The substitution $z=x/\sqrt{4t}$ and the one-dimensional Gaussian
integral give
$$\int_{\R^n}\Phi(x,t)\,dx
=\pi^{-n/2}\prod_{i=1}^n\int_{-\infty}^\infty e^{-z_i^2}\,dz_i=1.$$
\end{proof}

\begin{lemma}[The fundamental solution solves the heat equation]
\label{lem:heat-kernel-solves-heat-equation}
\lean{EvansLib.heatKernelSpatial_solves_heat}
\uses{def:fundamental-solution-heat-equation}
\leanok
\group{Heat Equation}



For $t>0$ and $x\in\R^n$,
$$\Phi_t(x,t)-\Delta_x\Phi(x,t)=0.$$
\end{lemma}
\begin{proof}
\leanok
\sketch Direct differentiation gives
$$\Phi_t=\Phi\left(-\frac n{2t}+\frac{|x|^2}{4t^2}\right),
\qquad
\Phi_{x_jx_j}=\Phi\left(\frac{x_j^2}{4t^2}-\frac1{2t}\right).$$
Summing over $j$ proves the identity.
\end{proof}

\subsubsection{Initial-value problem}
\label{sec:heat-initial-value-problem}

For initial data $g$, define
$$(9)\qquad
u(x,t):=\int_{\R^n}\Phi(x-y,t)g(y)\,dy
\qquad(t>0).$$

\begin{theorem}[Solution of initial-value problem]
\label{thm:heat-equation-initial-value-solution}
\uses{def:fundamental-solution-heat-equation, lem:integral-fundamental-solution-heat, lem:heat-convolution-smooth, lem:heat-convolution-solves-heat, lem:heat-convolution-attains-initial-data}
\dcref{ch2:2.3-thm-1}
\group{Heat Equation}

If $g\in C(\R^n)\cap L^\infty(\R^n)$ and $u$ is defined by (9), then
$u\in C^\infty(\R^n\times(0,\infty))$, $u_t-\Delta u=0$, and
$$\lim_{\substack{(x,t)\to(x^0,0)\\t>0}}u(x,t)=g(x^0)
\qquad(x^0\in\R^n).$$
\end{theorem}
\begin{proof}
\sketch On every strip $t\ge\delta>0$, all derivatives of the Gaussian are
bounded by integrable Gaussian multiples, so differentiation under the
integral gives smoothness and
$$u_t-\Delta u=\int_{\R^n}(\Phi_t-\Delta\Phi)(x-y,t)g(y)\,dy=0.$$
For the initial limit, write
$u(x,t)-g(x^0)=\int\Phi(x-y,t)[g(y)-g(x^0)]\,dy$ and split at
$B(x^0,\delta)$. Continuity bounds the near part; boundedness of $g$ and the
vanishing Gaussian tail as $t\downarrow0$ control the far part, uniformly for
$x\to x^0$.
\end{proof}

\begin{definition}[Heat convolution solution]
\label{def:heat-convolution-solution}
\lean{EvansLib.heatSolution}
\uses{def:fundamental-solution-heat-equation}
\leanok
\group{Heat Equation}



For $g:\R^n\to\R$, the \emph{heat convolution solution} is
$$u(x,t):=\int_{\R^n}\Phi(x-y,t)g(y)\,dy
\qquad(x\in\R^n,\ t>0).$$
\end{definition}

\begin{lemma}[The heat convolution is smooth]
\label{lem:heat-convolution-smooth}
\lean{EvansLib.heatSolution_contDiffOn}
\uses{def:heat-convolution-solution, def:fundamental-solution-heat-equation}
\leanok
\group{Heat Equation}



If $g\in C(\R^n)$ has compact support, its heat convolution is
$C^\infty$ on $\R^n\times(0,\infty)$.
\end{lemma}
\begin{proof}
\leanok
\sketch Near a fixed $(x_0,t_0)$, insert a smooth compactly supported cutoff
in the kernel that equals $1$ on all differences $x-y$ with $x$ near $x_0$
and $y\in\spt g$. The resulting compactly supported smooth kernel has a smooth
parametric convolution, which locally equals $u$.
\end{proof}

\begin{lemma}[The heat convolution solves the heat equation]
\label{lem:heat-convolution-solves-heat}
\lean{EvansLib.heatSolution_solves_heat}
\uses{def:heat-convolution-solution, lem:heat-kernel-solves-heat-equation}
\leanok
\group{Heat Equation}



If $g\in C(\R^n)$ has compact support, then its heat convolution satisfies
$$\partial_tu(x,t)=\sum_{j=1}^n\partial_{x_j}^2u(x,t)$$
for every $x\in\R^n$ and $t>0$.
\end{lemma}
\begin{proof}
\leanok
\sketch Compact support supplies a common integrable bound for the required
parameter derivatives. Differentiate under the integral and use
$\Phi_t=\Delta_x\Phi$.
\end{proof}

\begin{remark}[Dirac measure notation; infinite propagation speed]
\label{rem:heat-equation-infinite-propagation-speed}
\uses{thm:heat-equation-initial-value-solution}
\dcref{ch2:2.3.1-rem-1}
\group{Heat Equation}

The heat kernel satisfies symbolically
$$\Phi_t-\Delta\Phi=0\ \text{for }t>0,
\qquad
\Phi(\cdot,0)=\delta_0.$$
If $g$ is bounded, continuous, nonnegative, and not identically zero, then its
heat convolution is positive at every $x\in\R^n$ and every $t>0$.
\end{remark}

\begin{lemma}[Infinite propagation speed]
\label{lem:heat-infinite-propagation-speed}
\lean{EvansLib.heatSolution_pos_of_bounded}
\uses{def:heat-convolution-solution}
\group{Heat Equation}



If $g\in C(\R^n)$ has compact support, $g\ge0$, and $g$ is positive somewhere,
then its heat convolution satisfies $u(x,t)>0$ for all $x\in\R^n$ and $t>0$.
\end{lemma}
\begin{proof}
\leanok
\sketch The Gaussian is everywhere positive for $t>0$, while continuity makes
$g$ positive on a nonempty open set. The nonnegative integrand is therefore
positive on a set of positive measure.
\end{proof}

\begin{lemma}[The heat convolution attains its initial data]
\label{lem:heat-convolution-attains-initial-data}
\lean{EvansLib.heatSolution_tendsto_initial}
\uses{def:heat-convolution-solution, lem:integral-fundamental-solution-heat}
\leanok
\group{Heat Equation}



If $g\in C(\R^n)$ has compact support, then
$$\lim_{\substack{(x,t)\to(x^0,0)\\t>0}}u(x,t)=g(x^0).$$
\end{lemma}
\begin{proof}
\leanok
\sketch Using $\int\Phi(\cdot,t)=1$, write
$$u(x,t)-g(x)=\int_{\R^n}\Phi(z,t)[g(x-z)-g(x)]\,dz.$$
Uniform continuity controls $|z|<\delta$; boundedness of $g$ controls the
complement by the Gaussian mass outside $B(0,\delta)$, which tends to zero
after the scaling $z=\sqrt t\,w$. Finally let $x\to x^0$.
\end{proof}

\subsubsection{Nonhomogeneous problem}
\label{sec:heat-nonhomogeneous-problem}

For zero initial data, Duhamel's formula is
$$(13)\qquad
u(x,t):=\int_0^t\int_{\R^n}
\Phi(x-y,t-s)f(y,s)\,dy\,ds.$$

\begin{theorem}[Solution of nonhomogeneous problem]
\label{thm:heat-equation-nonhomogeneous-solution}
\uses{def:fundamental-solution-heat-equation}
\dcref{ch2:2.3-thm-2}
\group{Heat Equation}

Assume $f\in C_1^2(\R^n\times[0,\infty))$ has compact support and define
$u$ by (13). Then $u\in C_1^2(\R^n\times(0,\infty))$,
$$u_t-\Delta u=f,$$
and $u(x,t)\to0$ as $(x,t)\to(x^0,0)$ with $t>0$.
\end{theorem}
\begin{proof}
\sketch Rewrite
$$u(x,t)=\int_0^t\int_{\R^n}\Phi(y,s)f(x-y,t-s)\,dy\,ds.$$
Differentiating the smooth compactly supported factor gives the asserted
regularity. Split the heat-operator calculation at $s=\varepsilon$. The part
$0<s<\varepsilon$ is $O(\varepsilon)$ because $\int\Phi=1$. Integrating by
parts on $\varepsilon<s<t$ moves $-\partial_s-\Delta_y$ from $f$ to $\Phi$;
the interior term vanishes and the endpoint at $s=\varepsilon$ tends to
$f(x,t)$. Finally
$\|u(\cdot,t)\|_\infty\le t\|f\|_\infty\to0$.
\end{proof}

\begin{remark}[Combined Duhamel formula]
\label{rem:heat-equation-duhamel-combined}
\uses{thm:heat-equation-initial-value-solution, thm:heat-equation-nonhomogeneous-solution}
\dcref{ch2:2.3.1-rem-2}
\group{Heat Equation}

Under the preceding hypotheses,
$$(17)\qquad
u(x,t)=\int_{\R^n}\Phi(x-y,t)g(y)\,dy
+\int_0^t\int_{\R^n}\Phi(x-y,t-s)f(y,s)\,dy\,ds$$
solves $u_t-\Delta u=f$ with $u(\cdot,0)=g$.
\end{remark}

\subsection{Mean-Value Formula}
\label{sec:heat-mean-value-formula}

\begin{definition}[Parabolic cylinder]
\label{def:parabolic-cylinder-heat-ball}
\lean{EvansLib.parabolicCylinder}
\leanok
\dcref{ch2:2.3.2-def-1}
\group{Heat Equation}


For a bounded open $U\subset\R^n$ and $T>0$, define
$$U_T:=U\times(0,T].$$
\end{definition}

\begin{definition}[Parabolic boundary]
\label{def:parabolic-boundary-heat-ball}
\lean{EvansLib.parabolicBoundary}
\uses{def:parabolic-cylinder-heat-ball}
\leanok
\dcref{ch2:2.3.2-def-2}
\group{Heat Equation}



The parabolic boundary is
$$\Gamma_T:=\overline U_T\setminus U_T.$$
It consists of $U\times\{0\}$ and $\partial U\times[0,T]$, but not the top
interior $U\times\{T\}$.
\end{definition}

\begin{definition}[Heat ball]
\label{def:heat-ball}
\lean{EvansLib.heatBall}
\uses{def:fundamental-solution-heat-equation}
\leanok
\dcref{ch2:2.3.2-def-3}
\group{Heat Equation}



For $x\in\R^n$, $t\in\R$, and $r>0$, define
$$E(x,t;r):=
\left\{(y,s):s\le t,\ \Phi(x-y,t-s)\ge r^{-n}\right\}.$$
\end{definition}

\begin{theorem}[Mean-value property for the heat equation]
\label{thm:mean-value-property-heat-equation}
\uses{def:fundamental-solution-heat-equation, def:parabolic-cylinder-heat-ball, def:heat-ball}
\dcref{ch2:2.3-thm-3}
\group{Heat Equation}

If $u\in C_1^2(U_T)$ solves the heat equation, then
$$(19)\qquad
u(x,t)=\frac1{4r^n}\iint_{E(x,t;r)}
u(y,s)\frac{|x-y|^2}{(t-s)^2}\,dy\,ds$$
whenever $E(x,t;r)\subset U_T$.
\end{theorem}
\begin{proof}
\sketch Translate $(x,t)$ to $(0,0)$, write $E(r)=E(0,0;r)$, and set
$$\phi(r)=\frac1{r^n}\iint_{E(r)}u(y,s)\frac{|y|^2}{s^2}\,dy\,ds.$$
After scaling to $E(1)$, differentiate in $r$. With
$$\psi=-\frac n2\log(-4\pi s)+\frac{|y|^2}{4s}+n\log r,$$
one has $\psi=0$ on $\partial E(r)$. Spatial and time integration by parts,
together with $u_s-\Delta u=0$, shows $\phi'(r)=0$. The scaling identity
$$r^{-n}\iint_{E(r)}\frac{|y|^2}{s^2}\,dy\,ds=4$$
then gives $\phi(r)=4u(0,0)$.
\end{proof}

\subsection{Properties of Solutions}
\label{sec:heat-properties-of-solutions}

\subsubsection{Strong maximum principle, uniqueness}
\label{sec:heat-strong-maximum-principle}

\begin{theorem}[Strong maximum principle for the heat equation]
\label{thm:strong-maximum-principle-heat-equation}
\uses{thm:mean-value-property-heat-equation}
\dcref{ch2:2.3-thm-4}
\group{Maximum Principles}

Let $u\in C_1^2(U_T)\cap C(\overline U_T)$ solve the heat equation.

(i) $\max_{\overline U_T}u=\max_{\Gamma_T}u$.

(ii) If $U$ is connected and $u$ attains its maximum at
$(x_0,t_0)\in U_T$, then $u$ is constant on $U\times[0,t_0]$.
\end{theorem}
\begin{proof}
\sketch At an interior maximum, the positive-weight heat-ball mean-value
formula forces $u$ to equal its maximum throughout every sufficiently small
backward heat ball. Propagate this equality backward along any decreasing-time
line segment. Polygonal connectivity of $U$ and a decreasing sequence of times
then connect $(x_0,t_0)$ to every $(x,t)$ with $t<t_0$. This proves (ii), and
(i) follows unless the solution is constant.
\end{proof}

\begin{remark}[Time-directedness of the strong maximum principle]
\label{rem:strong-maximum-principle-heat-time-intuition}
\uses{thm:strong-maximum-principle-heat-equation}
\dcref{ch2:2.3.3-rem-4}
\group{Maximum Principles}

An interior maximum at time $t_0$ forces constancy only at times not exceeding
$t_0$; later boundary data may change the solution.
\end{remark}

\begin{remark}[Strict positivity from boundary and initial data]
\label{rem:strong-maximum-principle-heat-positivity}
\uses{thm:strong-maximum-principle-heat-equation}
\dcref{ch2:2.3.3-rem-5}
\group{Maximum Principles}

Let $U$ be connected. If $u\in C_1^2(U_T)\cap C(\overline U_T)$ solves the
heat equation, vanishes on $\partial U\times[0,T]$, and has nonnegative initial
data positive somewhere, then $u>0$ throughout $U_T$.
\end{remark}

\begin{theorem}[Weak maximum principle for the heat equation]
\label{thm:weak-maximum-principle-heat}
\lean{EvansLib.exists_parabolicBoundary_isMaxOn}
\uses{def:parabolic-cylinder-heat-ball, def:parabolic-boundary-heat-ball}
\leanok
\group{Maximum Principles}



Let $U\subset\R^n$ be open, bounded, and nonempty, and let $T>0$. If
$u\in C(\overline U_T)$ is $C^2$ and satisfies $u_t=\Delta u$ on
$U\times(0,T)$, then
$$\max_{\overline U_T}u=\max_{\Gamma_T}u.$$
\end{theorem}
\begin{proof}
\leanok
\sketch For $T'<T$ and $\varepsilon>0$, set $v=u-\varepsilon t$ on
$\overline U_{T'}$. At an interior maximum, $v_t\ge0$ and $\Delta v\le0$, but
$$v_t=\Delta v-\varepsilon<0,$$
a contradiction. Thus the maximum lies on $\Gamma_{T'}\subset\Gamma_T$.
Let $\varepsilon\downarrow0$, then $T'\uparrow T$, using continuity at the top
slice.
\end{proof}

\begin{theorem}[Uniqueness on bounded domains]
\label{thm:uniqueness-heat-bounded-domain}
\lean{EvansLib.eqOn_closure_of_eqOn_parabolicBoundary}
\uses{thm:strong-maximum-principle-heat-equation, thm:weak-maximum-principle-heat}
\leanok
\dcref{ch2:2.3-thm-5}
\group{Maximum Principles}



For $g\in C(\Gamma_T)$ and $f\in C(U_T)$, there is at most one
$u\in C_1^2(U_T)\cap C(\overline U_T)$ satisfying
$$(22)\qquad
\begin{cases}
u_t-\Delta u=f&\text{in }U_T,\\
u=g&\text{on }\Gamma_T.
\end{cases}$$
\end{theorem}
\begin{proof}
\uses{thm:weak-maximum-principle-heat}
\leanok
\sketch The difference $w$ of two solutions satisfies the homogeneous heat
equation and vanishes on $\Gamma_T$. Apply the weak maximum principle to
$w$ and $-w$.
\end{proof}

\begin{definition}[Backward-Gaussian comparison kernel]
\label{def:comparison-kernel}
\lean{EvansLib.compKernelSpaceTime}
\uses{def:fundamental-solution-heat-equation}
\leanok
\group{Heat Equation}



For $y\in\R^n$, $\tau\in\R$, and $t<\tau$, define
$$K(x,t):=(\tau-t)^{-n/2}
e^{|x-y|^2/[4(\tau-t)]}.$$
\end{definition}

\begin{lemma}[The comparison kernel solves the heat equation]
\label{lem:comparison-kernel-solves-heat}
\lean{EvansLib.compKernelSpaceTime_solvesHeat}
\uses{def:comparison-kernel}
\leanok
\group{Heat Equation}



For $t<\tau$,
$$K_t-\Delta_xK=0.$$
\end{lemma}
\begin{proof}
\leanok
\sketch Put $s=\tau-t$. Direct differentiation gives
$$\Delta_xK=K\left(\frac{|x-y|^2}{4s^2}+\frac n{2s}\right)
=-\partial_sK=K_t.$$
\end{proof}

\begin{lemma}[The comparison function solves the heat equation]
\label{lem:comparison-function-solves-heat}
\lean{EvansLib.sub_const_mul_compKernel_solvesHeat}
\uses{lem:comparison-kernel-solves-heat}
\leanok
\group{Heat Equation}



If $u$ solves the heat equation below time $\tau$ and $\mu\in\R$, then
$v=u-\mu K$ also solves it.
\end{lemma}
\begin{proof}
\leanok
\sketch Apply linearity of $\partial_t-\Delta$.
\end{proof}

\begin{theorem}[Maximum principle for the Cauchy problem]
\label{thm:maximum-principle-heat-cauchy-problem}
\lean{EvansLib.heat_cauchy_maxPrinciple}
\uses{thm:maximum-principle-heat-cauchy-problem-smalltime}
\leanok
\dcref{ch2:2.3-thm-6}
\group{Maximum Principles}



Suppose $u\in C_1^2(\R^n\times(0,T])\cap C(\R^n\times[0,T])$ solves
$$u_t-\Delta u=0,
\qquad
u(x,0)=g(x),$$
and for constants $A,a>0$ satisfies
$$(24)\qquad u(x,t)\le Ae^{a|x|^2}
\quad(x\in\R^n,\ 0\le t\le T).$$
Then
$$\sup_{\R^n\times[0,T]}u=\sup_{\R^n}g.$$
\end{theorem}
\begin{proof}
\sketch First suppose $4aT<1$ and choose $\varepsilon>0$ with
$4a(T+\varepsilon)<1$. For fixed $y$ and $\mu>0$, subtract the comparison
kernel with singular time $T+\varepsilon$:
$$v(x,t)=u(x,t)-\mu(T+\varepsilon-t)^{-n/2}
e^{|x-y|^2/[4(T+\varepsilon-t)]}.$$
On a ball cylinder $B(y,r)\times(0,T]$, $v$ solves the heat equation. Its
bottom value is at most $\sup g$, and the kernel dominates the growth bound
on the lateral boundary for all sufficiently large $r$. The bounded-domain
maximum principle gives $v(y,t)\le\sup g$; let $\mu\downarrow0$. For general
$T$, repeat the small-time result on successive slabs of length $1/(8a)$.
The reverse inequality holds on the initial slice.
\end{proof}

\begin{theorem}[Maximum principle for the Cauchy problem, small-time case]
\label{thm:maximum-principle-heat-cauchy-problem-smalltime}
\lean{EvansLib.heat_cauchy_maxPrinciple_smallTime}
\uses{lem:comparison-function-solves-heat, thm:weak-maximum-principle-heat}
\group{Maximum Principles}



Let $u\in C(\R^n\times[0,T])$ be $C^2$ and solve the heat equation on
$\R^n\times(0,T)$. Suppose $u(x,0)\le M$,
$u(x,t)\le Ae^{a|x|^2}$, and $4aT<1$. Then
$$u\le M\quad\text{on }\R^n\times[0,T].$$
\end{theorem}
\begin{proof}
\sketch Choose $\varepsilon>0$ with $4a(T+\varepsilon)<1$, compare $u$ with
$\mu K$ on $B(y,r)\times(0,T]$, and apply the weak maximum principle. The
initial boundary is bounded by $M$, while the backward Gaussian dominates the
growth term on the lateral boundary as $r\to\infty$. Let $\mu\downarrow0$.
\end{proof}

\begin{theorem}[Uniqueness for Cauchy problem]
\label{thm:uniqueness-heat-cauchy-problem}
\lean{EvansLib.heat_cauchy_uniqueness}
\uses{thm:maximum-principle-heat-cauchy-problem}
\leanok
\dcref{ch2:2.3-thm-7}
\group{Maximum Principles}



For $g\in C(\R^n)$ and $f\in C(\R^n\times[0,T])$, there is at most one
$u\in C_1^2(\R^n\times(0,T])\cap C(\R^n\times[0,T])$ satisfying
$$u_t-\Delta u=f,
\qquad
u(x,0)=g(x),$$
and
$$|u(x,t)|\le Ae^{a|x|^2}$$
for constants $A,a>0$.
\end{theorem}
\begin{proof}
\sketch Apply the Cauchy maximum principle to both signs of the difference of
two solutions.
\end{proof}

\begin{theorem}[Uniqueness for Cauchy problem, small-time case]
\label{thm:uniqueness-heat-cauchy-problem-smalltime}
\lean{EvansLib.heat_cauchy_uniqueness_smallTime}
\uses{thm:maximum-principle-heat-cauchy-problem-smalltime}
\group{Maximum Principles}



Let $u,v\in C(\R^n\times[0,T])$ be $C^2$ on $\R^n\times(0,T)$, agree at
$t=0$, and suppose $u-v$ solves the homogeneous heat equation. If
$|u|,|v|\le Ae^{a|x|^2}$ and $4aT<1$, then $u=v$.
\end{theorem}
\begin{proof}
\sketch Apply the small-time maximum principle to $\pm(u-v)$ with growth
constant $2A$ and initial bound $0$.
\end{proof}

\begin{remark}[Nonuniqueness without a growth bound]
\label{rem:heat-equation-nonphysical-solutions}
\uses{thm:uniqueness-heat-cauchy-problem}
\dcref{ch2:2.3.3-rem-7}
\group{Maximum Principles}

Without a growth restriction, the zero-data Cauchy problem
$$u_t-\Delta u=0,
\qquad
u(x,0)=0$$
has infinitely many solutions; all nonzero examples grow rapidly as
$|x|\to\infty$ \cite{John1982}.
\end{remark}

\subsubsection{Regularity}
\label{sec:heat-regularity}

\begin{theorem}[Smoothness]
\label{thm:heat-equation-smoothness}
\uses{thm:heat-equation-nonhomogeneous-solution, thm:uniqueness-heat-cauchy-problem}
\dcref{ch2:2.3-thm-8}
\group{Heat Equation}

Every $u\in C_1^2(U_T)$ solving the heat equation belongs to
$C^\infty(U_T)$.
\end{theorem}
\begin{proof}
\sketch Fix a closed parabolic cylinder $C=C(x_0,t_0;r)\subset U_T$ and a
cutoff $\zeta$ equal to $1$ on a smaller cylinder $C'$ and vanishing near the
parabolic boundary of $C$. For $v=\zeta u$,
$$v_t-\Delta v=(\zeta_t-\Delta\zeta)u-2D\zeta\cdot Du=:\widetilde f.$$
Duhamel's formula and Cauchy uniqueness represent $v$ by the heat convolution
of $\widetilde f$. Integrating the $Du$ term by parts gives, on a still smaller
cylinder $C''$,
$$u(x,t)=\iint_CK(x,t,y,s)u(y,s)\,dy\,ds,$$
where $K$ is smooth because it vanishes near the heat-kernel singularity.
Mollification justifies the identity at the stated regularity. Differentiation
under this integral gives $u\in C^\infty(C'')$.
\end{proof}

\subsubsection{Local estimates for solutions of the heat equation}
\label{sec:heat-local-estimates}

\begin{theorem}[Estimates on derivatives]
\label{thm:heat-equation-derivative-estimates}
\uses{thm:heat-equation-smoothness}
\dcref{ch2:2.3-thm-9}
\group{Heat Equation}

For every $k,l\in\N$ there is $C_{kl}$ such that
$$\max_{C(x,t;r/2)}|D_x^kD_t^lu|
\le\frac{C_{kl}}{r^{k+2l+n+2}}
\|u\|_{L^1(C(x,t;r))}$$
whenever $C(x,t;r)\subset U_T$ and $u$ solves the heat equation.
\end{theorem}
\begin{proof}
\sketch On the unit cylinder, the smooth-kernel representation from the
regularity proof yields
$$|D_x^kD_t^lu|\le C_{kl}\|u\|_{L^1(C(1))}$$
on the half cylinder. For $C(r)$, apply this to $v(x,t)=u(rx,r^2t)$. Since
$D_x^kD_t^lv=r^{k+2l}(D_x^kD_t^lu)(rx,r^2t)$ and
$\|v\|_{L^1(C(1))}=r^{-n-2}\|u\|_{L^1(C(r))}$, the stated exponent follows.
\end{proof}

\begin{remark}[Spatial analyticity]
\label{rem:heat-equation-spatial-analyticity}
\uses{thm:heat-equation-derivative-estimates}
\dcref{ch2:2.3.3-rem-8}
\group{Heat Equation}

For each fixed $0<t\le T$, the map $x\mapsto u(x,t)$ is analytic, but
$t\mapsto u(x,t)$ need not be analytic \cite{Mikhailov1978}.
\end{remark}

\subsection{Energy Methods}
\label{sec:heat-energy-methods}

\subsubsection{Uniqueness}
\label{sec:heat-energy-uniqueness}

\begin{theorem}[Uniqueness, energy method]
\label{thm:uniqueness-heat-energy-method}
\lean{EvansLib.eqOn_closure_of_eqOn_parabolicBoundary_anisotropic}
\leanok
\dcref{ch2:2.3-thm-10}
\group{Energy Methods}
For $g\in C(\Gamma_T)$ and $f\in C(U_T)$, there is at most one
$u\in C_1^2(\overline U_T)$ satisfying
$$u_t-\Delta u=f\ \text{in }U_T,
\qquad
u=g\ \text{on }\Gamma_T.$$
\end{theorem}
\begin{proof}
  \uses{lem:uniqueness-heat-bounded-domain-anisotropic}
\sketch For the difference $w$ of two solutions, define
$e(t)=\int_Uw(x,t)^2\,dx$. Since $w=0$ on the parabolic boundary,
$$e'(t)=2\int_Uw\Delta w\,dx=-2\int_U|Dw|^2\,dx\le0.$$
As $e(0)=0$, one has $w=0$.
\end{proof}

\subsubsection{Backwards uniqueness}
\label{sec:heat-backwards-uniqueness}

\begin{theorem}[Backwards uniqueness]
\label{thm:backwards-uniqueness-heat-equation}
\uses{thm:uniqueness-heat-energy-method}
\dcref{ch2:2.3-thm-11}
\group{Energy Methods}

Let $u,\widetilde u\in C^2(\overline U_T)$ solve the heat equation with the
same data on $\partial U\times[0,T]$. If $u(x,T)=\widetilde u(x,T)$ on $U$,
then $u=\widetilde u$ throughout $U_T$.
\end{theorem}
\begin{proof}
\sketch Put $w=u-\widetilde u$ and $e(t)=\int_Uw^2$. Integration by parts and
$w_t=\Delta w$ give
$$e'=-2\int_U|Dw|^2,
\qquad
e''=4\int_U(\Delta w)^2,$$
and Cauchy--Schwarz yields
$$(e')^2\le ee''.$$
Where $e>0$, this says $(\log e)''\ge0$. If $e$ were positive before a first
later zero, log-convexity would give
$$e((1-\tau)t_1+\tau t_2)
\le e(t_1)^{1-\tau}e(t_2)^\tau=0,$$
a contradiction. Since $e(T)=0$, one has $e\equiv0$.
\end{proof}

\section{Wave Equation}
\label{sec:wave-equation}

Let $U\subset\R^n$ be open. The homogeneous and nonhomogeneous wave equations
are
$$(1)\qquad u_{tt}-\Delta u=0,
\qquad
(2)\qquad u_{tt}-\Delta u=f,$$
and $\Box u:=u_{tt}-\Delta u$.

\subsection{Solution by Spherical Means}
\label{sec:wave-solution-spherical-means}

\subsubsection{Solution for $n = 1$, d'Alembert's formula}
\label{sec:wave-dalembert-formula}

Consider
$$(3)\qquad
\begin{cases}
u_{tt}-u_{xx}=0&\text{in }\R\times(0,\infty),\\
u=g,\quad u_t=h&\text{on }\R\times\{t=0\}.
\end{cases}$$

\begin{definition}[d'Alembert's formula]
\label{def:dalembert-formula}
\lean{EvansLib.dAlembert}
\leanok
\group{Wave Equation}


For $g,h:\R\to\R$, d'Alembert's solution is
$$u(x,t):=\frac12[g(x+t)+g(x-t)]
+\frac12\int_{x-t}^{x+t}h(y)\,dy.$$
\end{definition}

\begin{lemma}[Second-order chain rule along a linear form]
\label{lem:second-order-chain-rule-linear-form}
\lean{EvansLib.iteratedFDeriv_two_comp_clm}
\leanok
\group{Wave Equation}


Let $E$ be a normed real vector space, $\varphi\in C^2(\R)$, and
$L:E\to\R$ a continuous linear form. For $p,v,w\in E$,
$$D^2(\varphi\circ L)(p)(v,w)=\varphi''(Lp)(Lv)(Lw).$$
\end{lemma}
\begin{proof}
\leanok
\sketch Apply the chain rule twice, using that the derivative of $L$ is $L$.
\end{proof}

\begin{theorem}[Solution of wave equation, $n = 1$]
\label{thm:wave-equation-dalembert-solution}
\lean{EvansLib.dAlembert_isClassicalSolution}
\uses{def:dalembert-formula, lem:second-order-chain-rule-linear-form}
\leanok
\dcref{ch2:2.4-thm-1}
\group{Wave Equation}



If $g\in C^2(\R)$ and $h\in C^1(\R)$, then d'Alembert's solution belongs to
$C^2(\R^2)$ and satisfies $u_{tt}-u_{xx}=0$.
\end{theorem}
\begin{proof}
\leanok
\sketch Let $H(y)=\int_0^yh$. With $L_\pm(x,t)=x\pm t$, the solution is a
linear combination of $g\circ L_\pm$ and $H\circ L_\pm$. For each such term,
the second-order chain rule and $(L_\pm e_t)^2=(L_\pm e_x)^2=1$ give equality
of the second time and space derivatives.
\end{proof}

\begin{theorem}[Initial conditions for d'Alembert's solution]
\label{thm:dalembert-initial-conditions}
\lean{EvansLib.dAlembert_init,EvansLib.dAlembert_init_time}
\uses{def:dalembert-formula}
\leanok
\group{Wave Equation}



Under the preceding hypotheses,
$$u(x,0)=g(x),
\qquad
u_t(x,0)=h(x).$$
\end{theorem}
\begin{proof}
\leanok
\sketch Substitute $t=0$ for the displacement. Differentiating the formula
gives
$$u_t=\frac12[g'(x+t)-g'(x-t)]
+\frac12[h(x+t)+h(x-t)],$$
which equals $h(x)$ at $t=0$.
\end{proof}

\begin{theorem}[General solution of the 1D wave equation]
\label{thm:wave-1d-general-solution}
\lean{EvansLib.waveGeneralSol_isClassicalSolution}
\uses{lem:second-order-chain-rule-linear-form}
\leanok
\group{Wave Equation}



For $F,G\in C^2(\R)$,
$$u(x,t)=F(x+t)+G(x-t)$$
is a $C^2$ solution of $u_{tt}-u_{xx}=0$.
\end{theorem}
\begin{proof}
\leanok
\sketch Apply \cref{lem:second-order-chain-rule-linear-form} to
$F\circ L_+$ and $G\circ L_-$.
\end{proof}

\begin{remark}[General solution of the 1D wave equation]
\label{rem:wave-equation-general-solution-1d}
\uses{thm:wave-equation-dalembert-solution, thm:wave-1d-general-solution}
\dcref{ch2:2.4.1-rem-1}
\group{Wave Equation}

Every d'Alembert solution has the form $F(x+t)+G(x-t)$, reflecting the
factorization
$$\Box=(\partial_t+\partial_x)(\partial_t-\partial_x).$$
If $g\in C^k$ and $h\in C^{k-1}$, the solution is $C^k$ but need not be
smoother.
\end{remark}

For the half-line problem with $u(0,t)=0$, extend $g$ and $h$ oddly to
$\widetilde g$ and $\widetilde h$. D'Alembert's formula yields, for $x,t\ge0$,
$$(10)\qquad
u(x,t)=
\begin{cases}
\dfrac12[g(x+t)+g(x-t)]+\dfrac12\displaystyle\int_{x-t}^{x+t}h(y)\,dy
&x\ge t,\\[8pt]
\dfrac12[g(x+t)-g(t-x)]+\dfrac12\displaystyle\int_{t-x}^{t+x}h(y)\,dy
&x\le t.
\end{cases}$$

\subsubsection{Spherical means}
\label{sec:wave-spherical-means}

For $x\in\R^n$ and $r,t>0$, define the spherical averages
$$(12)\qquad U(x;r,t):=\fint_{\partial B(x,r)}u(y,t)\,dS(y),$$
and similarly
$$(13)\qquad
G(x;r):=\fint_{\partial B(x,r)}g(y)\,dS(y),
\qquad
H(x;r):=\fint_{\partial B(x,r)}h(y)\,dS(y).$$

\begin{lemma}[Euler--Poisson--Darboux equation]
\label{lem:euler-poisson-darboux-equation}
\dcref{ch2:2.4-lem-1}
\group{Wave Equation}
Fix $x\in\R^n$. If $u\in C^m$ solves the wave initial-value problem, then
$U\in C^m(\R_+\times[0,\infty))$ and
$$(14)\qquad
\begin{cases}
U_{tt}-U_{rr}-\dfrac{n-1}{r}U_r=0,\\
U(r,0)=G(r),\quad U_t(r,0)=H(r).
\end{cases}$$
\end{lemma}
\begin{proof}
  \uses{lem:wave-spherical-mean-smoothness, lem:wave-spherical-mean-time-derivative, lem:weak-divergence-form-radial-equation}
\sketch Differentiating spherical averages and applying the divergence theorem
gives
$$U_r=\frac r n\fint_{B(x,r)}\Delta u\,dy.$$
Since $\Delta u=u_{tt}$,
$$\partial_r(r^{n-1}U_r)=r^{n-1}U_{tt},$$
which is the Euler--Poisson--Darboux equation. Repeated differentiation of the
average gives the asserted regularity and the initial conditions follow from
those of $u$.
\end{proof}

\subsubsection{Solution for $n = 3, 2$, Kirchhoff's and Poisson's formulas}
\label{sec:wave-kirchhoff-poisson-formulas}

For $n=3$, $\widetilde U=rU$ satisfies the one-dimensional wave equation with
zero boundary value at $r=0$. Formula (10) gives Kirchhoff's formula
$$(22)\qquad
u(x,t)=\fint_{\partial B(x,t)}
\bigl[th(y)+g(y)+Dg(y)\cdot(y-x)\bigr],dS(y).$$

For $n=2$, regard the data as independent of a third coordinate and apply the
three-dimensional formula. Integrating over the two hemispheres gives
$$(27)\qquad
u(x,t)=\frac12\fint_{B(x,t)}
\frac{tg(y)+t^2h(y)+tDg(y)\cdot(y-x)}
{(t^2-|y-x|^2)^{1/2}}\,dy.$$

\subsubsection{Solution for odd $n$}
\label{sec:wave-solution-odd-n}

\begin{lemma}[Some useful identities]
\label{lem:wave-equation-useful-identities}
\lean{EvansLib.waveRadial_second_deriv,EvansLib.waveRadial_expansion, EvansLib.evansRadialCoeff_zero_eq_doubleFactorial, EvansLib.waveRadial_transform_of_epd}
\dcref{ch2:2.4-lem-2}
\group{Wave Equation}
Let $\phi\in C^{k+1}(\R)$ and $k\ge1$. Then
$$\frac{d^2}{dr^2}\left(\frac1r\frac d{dr}\right)^{k-1}
(r^{2k-1}\phi)
=\left(\frac1r\frac d{dr}\right)^k(r^{2k}\phi'),$$
and
$$\left(\frac1r\frac d{dr}\right)^{k-1}(r^{2k-1}\phi)
=\sum_{j=0}^{k-1}\beta_j^k r^{j+1}\phi^{(j)}(r),$$
where the constants $\beta_j^k$ are independent of $\phi$ and
$$\beta_0^k=1\cdot3\cdot5\cdots(2k-1).$$
\end{lemma}

Let $n=2k+1\ge3$ and define
$$\widetilde U=left(\frac1r\partial_r\right)^{k-1}(r^{2k-1}U),
\quad
\widetilde G=left(\frac1r\partial_r\right)^{k-1}(r^{2k-1}G),
\quad
\widetilde H=left(\frac1r\partial_r\right)^{k-1}(r^{2k-1}H).$$

\begin{lemma}[$\tilde{U}$ solves the one-dimensional wave equation]
\label{lem:transformed-spherical-mean-wave-equation}
\uses{lem:euler-poisson-darboux-equation, lem:wave-equation-useful-identities}
\dcref{ch2:2.4-lem-3}
\group{Wave Equation}

The transformed average satisfies
$$\widetilde U_{tt}-\widetilde U_{rr}=0,
\qquad
\widetilde U(r,0)=\widetilde G(r),
\qquad
\widetilde U_t(r,0)=\widetilde H(r),
\qquad
\widetilde U(0,t)=0.$$
\end{lemma}
\begin{proof}
\sketch The identities in \cref{lem:wave-equation-useful-identities} and
$n=2k+1$ give
$$\widetilde U_{rr}
=\left(\frac1r\partial_r\right)^{k-1}
\left[r^{2k-1}\left(U_{rr}+\frac{n-1}{r}U_r\right)\right]
=\widetilde U_{tt}.$$
The expansion in powers of $r$ gives the boundary value at $r=0$.
\end{proof}

With $\gamma_n=1\cdot3\cdot5\cdots(n-2)$, d'Alembert's formula yields
$$(31)\qquad
u(x,t)=\frac1{\gamma_n}\left[
\partial_t\left(\frac1t\partial_t\right)^{(n-3)/2}
\left(t^{n-2}\fint_{\partial B(x,t)}g\,dS\right)
+\left(\frac1t\partial_t\right)^{(n-3)/2}
\left(t^{n-2}\fint_{\partial B(x,t)}h\,dS\right)
\right].$$

\begin{theorem}[Solution of wave equation in odd dimensions]
\label{thm:wave-equation-odd-dimensions-solution}
\lean{EvansLib.oddWaveSolution_isSolutionOfIVP_of_order}
\uses{lem:transformed-spherical-mean-wave-equation}
\leanok
\dcref{ch2:2.4-thm-2}
\group{Wave Equation}

Let $n\ge3$ be odd, $m=(n+1)/2$, $g\in C^{m+1}(\R^n)$, and
$h\in C^m(\R^n)$. The function (31) belongs to
$C^2(\R^n\times[0,\infty))$, solves $u_{tt}-\Delta u=0$, and satisfies
$$u(x,t)\to g(x^0),
\qquad
u_t(x,t)\to h(x^0)$$
as $(x,t)\to(x^0,0)$ with $t>0$.
\end{theorem}
\begin{proof}
\sketch By linearity treat the displacement and velocity data separately.
For $g=0$, differentiate (31) and use
$$\partial_tH(x;t)=\frac t n\fint_{B(x,t)}\Delta h\,dy,
\qquad
\Delta_xH(x;t)=\fint_{\partial B(x,t)}\Delta h\,dS,$$
to obtain $u_{tt}=\Delta u$; the $h=0$ calculation is analogous. The expansion
in \cref{lem:wave-equation-useful-identities} gives the initial limits.
\end{proof}

\begin{remark}[Domain of dependence, regularity, and Huygens' principle]
\label{rem:wave-equation-odd-dimension-remarks}
\uses{thm:wave-equation-odd-dimensions-solution, thm:wave-equation-dalembert-solution}
\dcref{ch2:2.4.1-rem-2}
\group{Wave Equation}

For odd $n\ge3$, $u(x,t)$ depends only on the data and their required
derivatives on $\partial B(x,t)$. Unlike the one-dimensional formula, (31)
uses derivatives of $g$, so wave evolution does not improve regularity. Both
formulas exhibit finite propagation speed.
\end{remark}

\subsubsection{Solution for even $n$}
\label{sec:wave-solution-even-n}

For even $n$, extend $g,h$ to $\R^{n+1}$ independently of the last coordinate
and apply the odd-dimensional formula. Integrating over the two hemispheres
gives, with $\gamma_n=2\cdot4\cdots(n-2)n$,
$$(38)\qquad
u(x,t)=\frac1{\gamma_n}\left[
\partial_t\left(\frac1t\partial_t\right)^{(n-2)/2}
\left(t^n\fint_{B(x,t)}
\frac{g(y)}{(t^2-|y-x|^2)^{1/2}}\,dy\right)
+\left(\frac1t\partial_t\right)^{(n-2)/2}
\left(t^n\fint_{B(x,t)}
\frac{h(y)}{(t^2-|y-x|^2)^{1/2}}\,dy\right)
\right].$$

\begin{theorem}[Solution of wave equation in even dimensions]
\label{thm:wave-equation-even-dimensions-solution}
\uses{thm:wave-equation-odd-dimensions-solution}
\dcref{ch2:2.4-thm-3}
\group{Wave Equation}

Let $n\ge2$ be even, $m=(n+2)/2$, $g\in C^{m+1}(\R^n)$, and
$h\in C^m(\R^n)$. The function (38) belongs to
$C^2(\R^n\times[0,\infty))$, solves $u_{tt}-\Delta u=0$, and satisfies
$$u(x,t)\to g(x^0),
\qquad
u_t(x,t)\to h(x^0)$$
as $(x,t)\to(x^0,0)$ with $t>0$.
\end{theorem}
\begin{proof}
\sketch Apply \cref{thm:wave-equation-odd-dimensions-solution} in dimension
$n+1$ to the extensions independent of the final coordinate. Parametrizing
the two hemispheres reduces the spherical averages to the weighted ball
integrals in (38), and restriction to the hyperplane recovers $u$.
\end{proof}

\begin{remark}[Huygens' principle]
\label{rem:huygens-principle}
\uses{thm:wave-equation-odd-dimensions-solution, thm:wave-equation-even-dimensions-solution}
\dcref{ch2:2.4.1-rem-3}
\group{Wave Equation}

For odd $n\ge3$, data at $x$ affect the solution only on the boundary
$|y-x|=t$ of the forward cone. For even $n$, formula (38) uses the entire ball
$B(y,t)$, so the data continue to influence the cone interior after the
leading front passes.
\end{remark}

\subsection{Nonhomogeneous Problem}
\label{sec:wave-nonhomogeneous-problem}

For each $s\ge0$, let $w(x,t;s)$ solve the homogeneous wave equation for
$t>s$ with
$$w(\cdot,s;s)=0,
\qquad
w_t(\cdot,s;s)=f(\cdot,s),$$
and define
$$(41)\qquad u(x,t):=\int_0^t w(x,t;s)\,ds.$$

\begin{theorem}[Solution of nonhomogeneous wave equation]
\label{thm:wave-equation-nonhomogeneous-duhamel}
\uses{thm:wave-equation-odd-dimensions-solution, thm:wave-equation-even-dimensions-solution}
\dcref{ch2:2.4-thm-4}
\group{Wave Equation}

Let $n\ge2$ and $f\in C^{[n/2]+1}(\R^n\times[0,\infty))$. The function
in (41) belongs to $C^2(\R^n\times[0,\infty))$, satisfies
$$u_{tt}-\Delta u=f,$$
and has $u(\cdot,0)=u_t(\cdot,0)=0$.
\end{theorem}
\begin{proof}
\sketch The odd- or even-dimensional representation theorem supplies the
required regularity of $w$. Leibniz' rule gives
$$u_t=\int_0^t w_t\,ds,
\qquad
u_{tt}=f(x,t)+\int_0^t w_{tt}\,ds,$$
while $\Delta u=\int_0^t\Delta w\,ds=\int_0^t w_{tt}\,ds$. The initial data
follow directly.
\end{proof}

\begin{remark}[General solution of the nonhomogeneous wave equation]
\label{rem:wave-equation-general-nonhomogeneous-solution}
\uses{thm:wave-equation-nonhomogeneous-duhamel}
\dcref{ch2:2.4.2-rem-1}
\group{Wave Equation}

For arbitrary initial displacement and velocity, add the corresponding
homogeneous solution to the Duhamel solution (41).
\end{remark}

\begin{example}[Nonhomogeneous wave equation for $n=1,3$]
\label{ex:wave-equation-nonhomogeneous-n1-n3}
\uses{thm:wave-equation-nonhomogeneous-duhamel, thm:wave-equation-dalembert-solution}
\dcref{ch2:2.4.2-ex-1}
\group{Wave Equation}

For $n=1$, zero initial data give
$$(43)\qquad
u(x,t)=\frac12\int_0^t\int_{x-s}^{x+s}f(y,t-s)\,dy\,ds.$$
For $n=3$, Kirchhoff's formula gives the retarded potential
$$(44)\qquad
u(x,t)=\frac1{4\pi}\int_{B(x,t)}
\frac{f(y,t-|y-x|)}{|y-x|}\,dy.$$
\end{example}

\subsection{Energy Methods}
\label{sec:wave-energy-methods}

\subsubsection{Uniqueness}
\label{sec:wave-energy-uniqueness}

Let $U\subset\R^n$ be bounded with smooth boundary. Consider
$$(45)\qquad
\begin{cases}
u_{tt}-\Delta u=f&\text{in }U_T,\\
u=g&\text{on }\Gamma_T,\\
u_t=h&\text{on }U\times\{t=0\}.
\end{cases}$$

\begin{theorem}[Uniqueness for wave equation]
\label{thm:uniqueness-wave-equation-energy-method}
\dcref{ch2:2.4-thm-5}
\group{Energy Methods}
There is at most one $u\in C^2(\overline U_T)$ satisfying (45).
\end{theorem}
\begin{proof}
\sketch For the difference $w$ of two solutions, define
$$e(t)=\frac12\int_U(w_t^2+|Dw|^2)\,dx.$$
The homogeneous wave equation and integration by parts give
$$e'(t)=\int_Uw_t(w_{tt}-\Delta w)\,dx=0,$$
with no boundary term because $w=w_t=0$ on the lateral boundary. Initial data
give $e(0)=0$, so $w_t=Dw=0$ and $w=0$.
\end{proof}

\subsubsection{Domain of dependence}
\label{sec:wave-domain-of-dependence}

Fix $x_0\in\R^n$, $t_0>0$, and let
$$C=\{(x,t):0\le t\le t_0,\ |x-x_0|\le t_0-t\}.$$

\begin{theorem}[Finite propagation speed]
\label{thm:wave-equation-finite-propagation-speed}
\dcref{ch2:2.4-thm-6}
\group{Energy Methods}
If $u\in C^2$ solves the homogeneous wave equation and
$u=u_t=0$ on $B(x_0,t_0)\times\{0\}$, then $u=0$ throughout $C$.
\end{theorem}
\begin{proof}
\sketch Define the local energy
$$e(t)=\frac12\int_{B(x_0,t_0-t)}(u_t^2+|Du|^2)\,dx.$$
Differentiation over the shrinking ball, integration by parts, and the wave
equation give
$$e'(t)=\int_{\partial B(x_0,t_0-t)}
\left(\frac{\partial u}{\partial\nu}u_t-\frac12u_t^2-\frac12|Du|^2\right)dS.$$
Since $|\partial_\nu u\,u_t|\le(u_t^2+|Du|^2)/2$, one has $e'\le0$.
The initial assumptions give $e(0)=0$, hence $u_t=Du=0$ and then $u=0$ in
the cone.
\end{proof}

\section{Transport Equation with Constant Reaction}
\label{sec:transport-equation-constant-reaction}

%ORIGINAL-NUMBER: 1
\begin{theorem}[Problem 1: first-order equation with constant reaction term]
\label{ex:transport-reaction-equation}
\lean{EvansLib.transportReactionSolution_char_deriv, EvansLib.transportReactionSolution_init}
\uses{thm:transport-ivp-solution, thm:transport-initial-condition}
\group{Transport Equation}



For constants $c\in\R$, $b\in\R^n$, the function
$$u(x,t)=e^{-ct}g(x-tb)$$
satisfies
$$u_t+b\cdot Du+cu=0,
\qquad
u(x,0)=g(x).$$
\end{theorem}
\begin{proof}
\leanok
\sketch The transport factor is constant along $(b,1)$, while the derivative
of $e^{-ct}$ contributes $-cu$. At $t=0$ the exponential equals $1$.
\end{proof}

\section{References}
\label{sec:chapter2-references}

Further sources for the material in this chapter include
\cite{FeynmanLeightonSands1966},
\cite{GilbargTrudinger1983},
\cite{Mikhailov1978},
\cite{John1982}, \cite{Friedman1964},
\cite{Watson1973}, \cite{John1982}, \cite{Mikhailov1978},
\cite{Payne1975}, \cite{Antman1980}, \cite{Folland1976}, and
\cite{Strauss1992}.
